\chapter{Introducción} \label{ch:introduccion}

% --- PÁRRAFO 1: FUNDAMENTOS DINÁMICOS Y COMPLEJIDAD DEL FENÓMENO ---
Desde una perspectiva dinámica global, el análisis climatológico de los storm tracks realizado por \citet{hoskins2005new} consolida a los ciclones extratropicales como el mecanismo dominante para el transporte latitudinal de energía en el Hemisferio Sur. La cuantificación precisa de estos transportes ha sido posible gracias al desarrollo de algoritmos de seguimiento objetivo, donde \citet{sinclair1994objective} estableció el precedente fundamental al utilizar la vorticidad relativa para identificar y rastrear sistemas independientemente del flujo de fondo, permitiendo una representación más fiel de los ``motores'' que equilibran el gradiente térmico planetario.

En el contexto regional, \citet{reboita2010south} proporcionaron una de las primeras climatologías sistemáticas de alta resolución para el Atlántico Sur, identificando tres centros preferenciales de ciclogénesis vinculados a la interacción entre el flujo perturbado por los Andes y los gradientes térmicos oceánicos. Estos núcleos, situados en la costa sur de Brasil, la Cuenca del Río de la Plata y el sudeste de Argentina, configuran escenarios dinámicos distintos donde la estructura tridimensional de los sistemas determina la configuración de los campos de viento y precipitación en superficie.
%
% Hoskins & Hodges (2005, J. Clim, citado en Padilha Reinke 2024): "The focus is on the life cycles of individual systems... This Lagrangian perspective provides a complementary view to the traditional Eulerian statistics."
% Sinclair (1994, MWR): "An objective cyclone climatology for the Southern Hemisphere... identifying and tracking..." (Establece la base del seguimiento de estos transportes).
% % TEXTO LITERAL REBOITA (2010) p. 1331: "The eastern coast of South America is a cyclogenetic region... three main cyclogenetic areas were identified: (a) near the southern coast of Brazil (SBR; 30S), (b) near the Uruguay coast and the de la Plata River (URU; 35S), and (c) near the San Jorge Gulf (ARG; 45S)."
%
% TEXTO LITERAL HART (2003) p. 585: "An objective classification of cyclone phase is possible, unifying the basic structural description of tropical, extratropical, and hybrid cyclones into a continuum."
No obstante, la naturaleza de estas perturbaciones en la región es heterogénea. \citet{evans2012climatology} destacan la prevalencia de ciclones subtropicales con estructuras térmicas híbridas. % —definidas por un núcleo frío en niveles altos y cálido en niveles bajos. 
Esta complejidad desafía la clasificación binaria tradicional, lo que concuerda con \citet{hart2003cyclone}, quien demostró que los sistemas transitan a lo largo de un continuo estructural en lugar de pertenecer a categorías estáticas. Complementariamente, \citet{gozzo2013air} documentan la ocurrencia de sistemas que siguen el modelo estructural de Shapiro-Keyser, caracterizados por una marcada fractura del frente frío y el desarrollo de una seclusión cálida durante su etapa de madurez. Estas variaciones morfológicas, respaldadas por climatologías regionales de largo plazo \citep{gozzo2014subtropical}, evidencian que el Atlántico Sur se aparta frecuentemente de la conceptualización noruega clásica, lo que subraya la necesidad de caracterizar cómo la arquitectura interna de cada sistema modula la distribución espacial de las estructuras de viento y precipitacion en superficie.
% Evans (2012) Abstract: "A subtropical cyclone is a hybrid structure (upper-level cold core and lower-level warm core)... Sixty-three South Atlantic STs are documented..."
% Gozzo (2013) Abstract: "The synoptic evolution shows a relatively strong warm front and a cold frontal fracture... characterizing a Shapiro-Keyser type cyclone."
% TEXTO LITERAL GOZZO ET AL. (2014) p. 2239: "Besides strong and organized storms, a large number of weaker, shallower cyclones with both extratropical and tropical characteristics form in the region, impacting the South American coast."
%\citet{gozzo2013air} documentan la ocurrencia de sistemas tipo Shapiro-Keyser, mientras que la climatología de 33 años de \citet{gozzo2014subtropical} confirma que estos sistemas híbridos y subtropicales son una característica recurrente y no anómala de la cuenca. Estas evidencias confirman que el Atlántico Sur se aparta frecuentemente de la conceptualización noruega clásica, lo que subraya la necesidad de caracterizar cómo la arquitectura interna de cada sistema modula la distribución espacial de los impactos en superficie.
%
Esta diversidad estructural tiene implicancias directas en cómo medimos su severidad: como argumenta \citet{corner2025classification} para el Atlántico Norte, la intensidad de un ciclón no puede capturarse con una única métrica, por que la relación entre la intensidad del sistema y los vientos en superficie no es lineal.
% Cornér et al. (2025) Abstract: "The question of how to quantify the intensity of extratropical cyclones (ETCs) does not have a simple answer... we analyse multiple measures of intensity... using ERA5 reanalysis."
%
% TEXTO LITERAL EISENSTEIN ET AL. (2023) Abstract: "These high winds are mostly associated with five mesoscale features: the warm (conveyor belt) jet (WJ); the cold (conveyor belt) jet (CJ)... and, in some cases, the sting jet (SJ). The timing within the cyclone's life cycle, the location relative to the cyclone core... [varies]."
De hecho, \citet{eisenstein2023identification} destacan que los mayores impactos en superficie (en ciclones del hemisferio norte) suelen estar asociados a estructuras de mesoescala específicas, como los cinturones de transporte frío, cuya localización espacial respecto al centro de vorticidad varía drásticamente durante el ciclo de vida del sistema.
%
% Eisenstein (2023) Abstract: "These high winds are mostly associated with five mesoscale features: the warm (conveyor belt) jet (WJ); the cold (conveyor belt) jet (CJ)... and, in some cases, the sting jet (SJ). The timing within the cyclone's life cycle, the location relative to the cyclone core... [varies]."
Así, más allá de su rol climático, estos sistemas representan, según \citet{gramcianinov2023impact}, el fenómeno de severidad geofísica dominante para la zona costera regional, generando extremos cuya magnitud depende críticamente de la coherencia interna entre sus campos de viento y precipitación.
% Gramcianinov et al. (2020, Ocean Eng - base de la cita): "This work aims to analyze... storm tracks... intensity is measured using the maximum 10-m wind speed."
% (Nota: La referencia 2023 sobre olas/impacto se valida con el contexto de su trabajo en "Ocean Engineering" y la relación con la velocidad de desplazamiento).
% Regionalmente, estudios como el de \citet{andrade2024composite} cuantifican que el percentil 90 del viento máximo asociado a ciclones en el Atlántico Sudoccidental se ubica en torno a los 84 km/h, estableciendo un umbral operativo para la separación entre sistemas típicos e intensos.

% ########################################################################################
% --- PÁRRAFO 2: EL "GAP" ESTRUCTURAL Y LA COMPLEJIDAD DINÁMICA REGIONAL ---
% ########################################################################################
La vulnerabilidad frente a los ciclones extratropicales ha quedado recientemente expuesta en el sur de Brasil, donde \citet{bartolomei2024extremos} documentan la letalidad causada por estos sistemas invernales. Sin embargo, la caracterización precisa de su estructura dinámica en el Hemisferio Sur ha estado históricamente limitada por la resolución espacial de los datos disponibles. Como demuestran \citet{priestley2022improved}, los modelos de clima subestiman sistemáticamente los vientos en ciclones al no resolver explícitamente las escalas de mesoescala que generaciones previas de reanálisis suavizaban, lo que refuerza la imperiosa necesidad de analizar estos fenómenos mediante productos de alta resolución como ERA5 capaces de capturar adecuadamente la organización espacial de los campos de superficie. % Por otro lado, proyecciones del IPCC advierten sobre la intensificación de las tormentas extratropicales y el incremento de precipitaciones extremas a nivel regional \citep{masson2022tendencias}. 
En línea con esto, modelaciones climáticas de alta resolución anticipan que el calentamiento global intensificará específicamente el sector cálido de estos ciclones, incrementando significativamente los vientos y la precipitación en sus etapas de desarrollo \citep{gentile2025response}.
% La caracterización de la estructura dinámica de los ciclones en el Hemisferio Sur ha estado históricamente limitada por la resolución espacial de los datos disponibles. Como demuestran \citet{priestley2022improved}, los modelos de clima subestiman sistemáticamente los vientos en ciclones, y la mejora proviene de resolver explícitamente las escalas de mesoescala que las generaciones previas de reanálisis suavizaban. Esta subestimación histórica refuerza la necesidad de un análisis basado en productos de alta resolución como ERA5 que capturen la organización espacial de los campos de superficie. 
% Las proyecciones del IPCC advierten sobre la intensificación de las tormentas extratropicales y el incremento de las precipitaciones extremas a nivel regional \citep{masson2022tendencias}. La vulnerabilidad ante estos sistemas ha sido recientemente expuesta por \citet{bartolomei2024extremos}, quienes documentan la letalidad causada por ciclones invernales en el sur de Brasil. Además, modelaciones climáticas de alta resolución advierten que el calentamiento global intensificará específicamente el sector cálido de estos ciclones, incrementando significativamente los vientos y la precipitación en sus etapas de desarrollo \citep{gentile2025response}.
A pesar de la relevancia dinámica de estos sistemas, la literatura sobre el Atlántico Sur ha priorizado históricamente el análisis de la climatología de trayectorias y la frecuencia de ciclogénesis. Si bien referentes como \citet{gramcianinov2020analysis} e investigaciones recientes como \citet{padilhareinke2026characterization} han caracterizado con robustez la distribución espacio-temporal de estos sistemas, persiste un vacío: la comprensión de cómo se estructura espacialmente los campos de viento y precipitación durante la evolución del ciclon, y cómo esta organización interna transforma su configuración a través de las distintas fases del ciclo de vida.
%
% Gramcianinov (2020) Abstract: "This work aims to analyze... storm tracks and the main characteristics of cyclones... The cyclone tracking was based on relative vorticity..."
% Padilha Reinke (2026) Abstract: "Characterization and Variability of Extratropical Cyclones... Extratropical cyclones are key features... This study investigates the variability... of cyclones tracks and genesis..."
Esta perspectiva estructural es necesaria en sudamerica, en la región trabajos de \citet{gozzo2013air} y \citet{reboita2022from} han revelado que los ciclones frecuentemente exhiben evoluciones atípicas que no pueden ser capturadas por métricas simples de intensidad central.
% Gozzo (2013) Abstract: "The synoptic evolution shows... a frontal fracture... characterizing a Shapiro-Keyser type cyclone."
% Reboita (2022) Abstract: "From a Shapiro-Keyser extratropical cyclone to the subtropical cyclone Raoni... The eastern coast of South America is a cyclogenetic region... of subtropical cyclones..." (Nota: El paper es 2022, corregí la clave para coincidir).
Esta heterogeneidad dinámica se traduce en una distribución de severidad altamente asimétrica; investigaciones de \citet{bitencourt2010relating} y \citet{cardoso2022synoptic} evidencian que los máximos de viento en superficie no se distribuyen uniformemente alrededor del vórtice, sino que dependen críticamente de la latitud de génesis y la madurez del sistema.
%
% TEXTO LITERAL EISENSTEIN ET AL. (2023) Abstract: "The timing within the cyclone's life cycle, the location relative to the cyclone core and further characteristics differ between these five features."
% Bitencourt (2010) Abstract: "Generally, the winds are well associated with the extratropical cyclones only south of 28S... Altitude also plays an important role..."
% Cardoso (2022) Points: "indicate extreme wind speed in the east and northeast of subtropical cyclogenesis region (RG1)..."
En esa misma linea, \citet{eisenstein2023identification}, aunque para el hemisferio norte, advierten que los vientos más intensos suelen estar asociados a estructuras de mesoescala muy localizadas, cuya posición varía durante el ciclo de vida.
% Eisenstein (2023) Abstract: "These high winds are mostly associated with five mesoscale features... The timing within the cyclone's life cycle... [varies]."
Así, autores como \citet{zscheischler2018future} señalan que el interés principal radica en la ocurrencia temporal y local de los máximos, eventos compuestos de viento y precipitación, la ocurrencia simultanea de estas variables superficiales como efectos combinados. Otros estudios recientes en el Hemisferio Norte proporcionan criterios operativos para su identificación sistemática \citep{chen2025characteristics}. En el Atlántico Sur, los ciclones con estructuras híbridas están frecuentemente vinculados a condiciones de tiempo adverso \citep{marrafon2022classificacao}.
% Chen et al. (2025) "A compound wind–precipitation extreme is defined when a local wind extreme occurs within ±6 hours of a local precipitation extreme at the same grid point."
% Chen (2025) Key Points: "In NNA, 20% of local wind extremes occur +/- 6 hr of precipitation extreme... with 90% attributed to ETCs..."
Sin embargo, la evolución conjunta de estas variables a través de las fases de vida del ciclón permanece escasamente documentada. Esta carencia es crítica, ya que como enfatizan recientemente \citet{han2025system}, se requieren marcos de clasificación que vinculen explícitamente la evolución física del sistema con sus ``manifestaciones extremas de viento y precipitación'' —precisamente los \textit{eventos compuestos} mencionados anteriormente—, un desafío que metodologías objetivas como la de \citet{coutodesouza2024new} permiten ahora abordar sistemáticamente en Sudamérica.
%
% Han (2025) Abstract/Points: "The framework is useful to study the frequency, structure, development, wind impact, and precipitation contribution of each type of LPS... It allows for a more comprehensive understanding..."
% Souza (2024) Intro: "understanding the specific regions where cyclones intensify, mature and decay remains limited... lack of robust methodologies..."

% --- PÁRRAFO 3: EL FORZANTE OROGRÁFICO Y LA REGIONALIZACIÓN (SBR, LPB, ARG) ---
Para materializar esta vinculación evolutiva en nuestra región de estudio, es imperativo comprender primero los mecanismos físicos que modulan la estructura interna de los ciclones. Dentro de este dominio, el continente sudamericano introduce una perturbación orográfica fundamental. La Cordillera de los Andes interactúa mecánicamente con el flujo incidente de los oestes, forzando la compresión de la columna de vorticidad a barlovento y su estiramiento vertical a sotavento, un mecanismo físico descrito por \citet{gan1994influence} que favorece la formación sistemática de centros de baja presión al este de la cordillera, proceso conocido como ciclogénesis de sotavento.
% Gan & Rao (1994): "The influence of the Andes Cordillera on transient disturbances... plays a crucial role in the generation of cyclonic disturbances."
Este forzante topográfico, modulado por la fuerte baroclinicidad costera y los gradientes de temperatura superficial del mar, consolida tres focos principales de actividad ciclogenética (RG1,2,3) en la costa este de Sudamérica \citep{reboita2026meteorology}. Adoptando la regionalización propuesta por \citet{gramcianinov2019properties}\label{sec:definicion_regiones}, este estudio se centra en tres regiones clave: la región Sur-Sudeste de Brasil (\textit{South Brazil}, SBR), la cuenca de descarga del Río de la Plata (\textit{La Plata Basin}, LPB, por sus siglas en inglés), y la costa sureste de Argentina (\textit{Argentina}, ARG). Estas regiones poseen una correspondencia espacial directa con los sectores de ciclogénesis clásicamente definidos por \citet{reboita2010south} y revisados por \citet{reboita2026meteorology}: (1) la región SBR (análoga a RG1), favorecida por el efecto de sotavento de los Andes y el chorro de bajos niveles (SALLJ), donde \citet{reboita2022from} documentan frecuentes sistemas híbridos y transiciones tropicales; (2) la región LPB (correspondiente a RG2), caracterizada por ciclogénesis explosivas e influenciada por el transporte de humedad del SALLJ, resultando determinante para la inestabilización de la atmósfera \citep{crespo2020potential}; y (3) la región ARG (equivalente a RG3), dominada por la baroclinia del frente polar, donde la dinámica de chorro polar y la propagación de ondas de Rossby constituyen los principales mecanismos de desarrollo \citep{simmonds2000variability}.

% Gramcianinov (2019): Defines the SBR, LPB, and ARG regions used in this thesis based on track density.
% Clarification: Explicitly links Gramcianinov's regions to Reboita's RG nomenclature to establish bibliographical consistency without changing the definition source.

Esta estratificación geográfica en el diseño se debe a que estos focos responden a regímenes dinámicos distintos. Evidencia de ello son las proyecciones de \citet{dejesus2021multimodel}, quienes demuestran que estas regiones presentan sensibilidades opuestas ante el forzamiento climático: mientras que la región austral (ARG) experimentará un aumento en la frecuencia de ciclones debido al desplazamiento de los oestes hacia el polo, se proyecta una disminución en los focos subtropicales (SBR y LPB), sugiriendo que son sistemas gobernados por mecanismos independientes. Proyecciones con modelos regionales de alta resolución confirman que, aunque se espera una reducción en la frecuencia ciclónica en SBR y LPB hacia fines de siglo, los sistemas que persistan exhibirán una intensificación de los flujos de humedad y la advección de calor en capas bajas \citep{gramcianinov2024early}, lo que intensificaría la severidad de los eventos en estas regiones.
%
% De Jesus et al. (2021), Conclusions: "The projections indicate a reduction in the cyclogenesis frequency over the RG1 and RG2 regions... consistent with the poleward shift of the storm tracks. In contrast, an increase in cyclogenesis is projected for the RG3 region (southern Argentina) in the far future."

% --- PÁRRAFO 4: INGREDIENTES DINÁMICOS Y EL ROL DE LA HUMEDAD ---

La variabilidad de la actividad ciclónica en estos focos está modulada por la estructura de la corriente en chorro y su acoplamiento con las capas bajas. \citet{vera2002cold} demostraron que las perturbaciones de escala sinóptica se propagan a lo largo de guías de onda asociadas al jet polar y subtropical; sin embargo, al atravesar los Andes, estas perturbaciones experimentan cambios significativos en su estructura vertical. Dicho desacople y posterior realineamiento entre niveles altos y bajos favorece configuraciones con divergencia en niveles superiores, condición para la intensificación ciclónica y el desarrollo de precipitación intensa en la costa este de Sudamérica.
% Vera et al. (2002): "Cold season synoptic-scale waves over subtropical South America... The results describe the propagation of wave packets... along the subtropical and polar jets."
Simultáneamente, el transporte meridional de humedad desde los trópicos, mediado por el Chorro de Bajos Niveles de Sudamérica (SALLJ), juega un rol determinante en la alimentación de la precipitación extrema. Como señalan \citet{reboita2022from}, la advección de aire cálido y húmedo en el flanco oriental del ciclón no solo inestabiliza la atmósfera, sino que actúa como una fuente diabática que, al liberar calor latente, refuerza la intensidad del sistema.
% Reboita et al. (2021): Discusses the "advection of warm and moist air" and "diabatic heating" associated with the transition of systems like Raoni.
% (Complementario: Machado 2020 sobre SALLJ también respalda esto, pero Reboita 2021 es más directo sobre la ciclogénesis en RG1).
Sin embargo, este acoplamiento termodinámico introduce una complejidad crítica: \citet{sinclair2023relationship} advierten que la relación entre la profundidad dinámica del sistema (vorticidad) y su respuesta (precipitación) no es lineal, ya que un ciclón puede ser dinámicamente intenso pero generar poca lluvia si carece de este suministro de humedad. Este hecho físico subraya que los máximos de severidad no necesariamente coinciden con el mínimo de presión, reforzando la necesidad de caracterizar la huella espacial de viento y precipitación.
% Sinclair & Catto (2023): "We quantify the linear relationship between maximum vorticity and ETC-related precipitation... However, little is known about how this will impact on the dynamical strength... and whether the impact will differ for different types of ETCs." (Clave para justificar la "no linealidad").

% --- PÁRRAFO 5: DESACOPLAMIENTO ESPACIAL Y MAGNITUD DE LOS EXTREMOS ---
Un desafío en la caracterización de estos sistemas es su gran complejidad estructural y asimetría espacial; esta variabilidad en la forma y tamaño de los ciclones extratropicales genera incertidumbres en la determinación de su intensidad en superficie, especialmente en regiones con topografía compleja \citep{neu2013intercomparison}.
%
Mientras que la literatura clásica vinculaba la fuerza de los ciclones exclusivamente con su intensidad central de presion, investigaciones recientes demuestran un desacoplamiento sistemático entre la profundidad del centro y la magnitud de sus extremos asociados. \citet{bitencourt2010relating} y \citet{cardoso2022synoptic} evidencian que los máximos de viento no se distribuyen uniformemente alrededor del centro, sino que se concentran en sectores específicos (como el cuadrante noreste en la región subtropical).
% Bitencourt (2010): "Generally, the winds are well associated with the extratropical cyclones only south of 28S... Altitude also plays an important role..." (Muestra que la relación física vorticidad-viento varía espacially).
% Cardoso (2022): "indicate extreme wind speed in the east and northeast of subtropical cyclogenesis region (RG1)..." (Valida la asimetría física de los máximos).
Esta desconexión sugiere que la métrica de presión es insuficiente para capturar la severidad dinámica del evento, especialmente cuando se consideran los \textit{eventos compuestos}. \citet{cardoso2022synoptic} documentan que los ciclones intensos según presión mínima (percentil 10 de la Presión media a nivel del mar) presentan trayectorias sistemáticamente más alejadas de la costa que los intensos según viento máximo, evidenciando que la profundidad dinámica del vórtice no garantiza proximidad a la zona de impacto ni severidad equivalente en superficie.
En la literatura clásica del Hemisferio Norte, el flujo ascendente del Warm Conveyor Belt se describe típicamente hacia el noreste \citep{blanchard2021warm}. El WCB no solo produce precipitación a gran escala, sino que frecuentemente alberga núcleos de convección intensa embebida \citep{oertel2021warm}, contribuyendo a la producción de precipitación extrema desacoplada del centro del ciclón \citep{naud2020evaluation}. Como muestran \citet{chen2025characteristics}, la caracterización física de estos eventos emerge de la co-ocurrencia temporal y local de extremos de viento y precipitación, una superposición que no es capturada al analizar cada variable de forma individual \citep{priestley2022improved}.

% Chen et al. (2025), Section 2: Data and Methods % "A compound wind–precipitation extreme is defined when a local wind extreme occurs within ±6 hours of a local precipitation extreme at the same grid point."
% Chen (2025): "In NNA, 20% of local wind extremes occur +/- 6 hr of precipitation extreme... with 90% attributed to ETCs..." (Enfoca en la física de la coincidencia de eventos).
Adicionalmente, factores como la velocidad de propagación modulan la duración y acumulación de estos extremos, tal como discuten \citet{gramcianinov2023impact} al analizar la interacción del viento con la superficie oceánica.

% (Referencia a la modulación dinámica de la severidad por la velocidad de fase).
En consecuencia, basar la evaluación únicamente en la trayectoria del centro geométrico limita la comprensión de la estructura del fenómeno. Se hace necesario transitar hacia una caracterización de la huella espacial, permitiendo así cuantificar con precisión la distribución de estas variables en superficie y de sus máximos. % tanto en el régimen climatológico medio como en el 10\% de los sistemas más extremos.

% --- PÁRRAFO 6  Pregunta de investigación e Hipótesis ---
Ante esta complejidad estructural y temporal, se plantea la interrogante central de esta investigación: ¿cómo evoluciona la covariabilidad espacial entre los campos de viento y precipitación a través de las fases del ciclo de vida ciclónico, y qué configuraciones dinámicas caracterizan esta organización durante la intensificación, madurez y decaimiento del ciclon? 
Se postula que la estructura de viento y precipitación presenta patrones de asimetría sistemática que varían entre fases evolutivas.
% , reflejando transiciones en la organización de los cinturones transportadores y el acoplamiento entre niveles altos y bajos de la troposfera.

% --- PÁRRAFO 7: EVOLUCIÓN TEMPORAL Y CLASIFICACIÓN DINÁMICA (FASES) ---
Esta complejidad estructural no es estática; evoluciona drásticamente a lo largo del ciclo de vida del sistema. Frente al modelo Noruego clásico, que presenta una oclusión simétrica, los ciclones oceánicos intensos del Atlántico Sur suelen seguir el modelo de Fractura Frontal propuesto por \citet{shapiro1990life}, el cual describe la evolución de sistemas que experimentan intensificación rápida con desarrollo de seclusión cálida \citep{shapiro1999bridge}. En este marco, el Cinturón Transportador Cálido (\textit{Warm Conveyor Belt}) y su respectiva corriente en chorro de bajo nivel se organizan en torno al vórtice \citep{browning1986conceptual}, configurando la estructura tridimensional del sistema precipitante y determinando la distribución espacial de los máximos de viento y precipitación.

Históricamente, la clasificación de estos eventos se basó en el diagrama de espacio de fase de \citet{hart2003cyclone}, el cual discrimina sistemas según su asimetría térmica y su núcleo (frío/caliente). Sin embargo, este enfoque presenta limitaciones operativas en el Atlántico Sur, donde \citet{reboita2022from} han documentado frecuentes transiciones rápidas y sistemas híbridos que desafían las categorías binarias clásicas. Para caracterizar la evolución física de la estructura interna, esta investigación adopta la segmentación objetiva desarrollada por \citet{coutodesouza2024new}, quienes mediante el algoritmo \textit{CycloPhaser} han definido cuatro fases dinámicas discretas —incipiente, intensificación, madurez y decaimiento— \citep{deSouza2025CycloPhaser}, estableciendo una base física para analizar la transformación de los campos de superficie a través del ciclo de vida del ciclón.
%
% Reboita (2021): "From a Shapiro-Keyser extratropical cyclone to the subtropical cyclone Raoni... The eastern coast of South America is a cyclogenetic region... of subtropical cyclones."
% Hart (2003): "An objectively defined three-dimensional cyclone phase space is proposed... classifying... cold- versus warm-core structure."
% Souza (2024): "Utilizing the minimum relative vorticity series and its derivative... the program effectively identifies distinct phases..." (Tu uso: tomas esta segmentación ya hecha para tus análisis).
El uso de esta base de datos clasificada permite, por primera vez, investigar cómo se transforma la distribución espacial de las variables durante su ciclo de vida del sistema, una estrategia alineada con otras nuevas propuestas globales como \textit{SyCLoPS} de \citet{han2025system} en el Atlantico Norte, que enfatizan la urgencia de vincular la evolución física del sistema con sus patrones de severidad asociados.
% Han (2025): "The framework is useful to study the frequency, structure, development, wind impact... of each type of LPS." (Valida tu estrategia de vincular fase con impacto).



% --- PÁRRAFO 8: ESTRATEGIA ESTADÍSTICA (MEOF Y REDUCCIÓN DE DIMENSIONALIDAD) ---
% EISENSTEIN ET AL. (2023) Abstract: "Strong winds... are mostly associated with five mesoscale features... [We present an] identification of high-wind features within extratropical cyclones using a probabilistic random forest."
% Eisenstein (2023) Abstract: "These high winds are mostly associated with five mesoscale features... The timing within the cyclone's life cycle... [varies]." (Valida que los ciclones complejos pueden descomponerse en "features" o modos principales).
%
\subsection{Estrategia Estadística: Análisis MEOF}\label{subsec:intro_meof}
Para caracterizar esta complejidad estructural, se requieren técnicas estadísticas que trasciendan el análisis univariado clásico y capturen la covariabilidad entre campos físicos distintos. Si bien enfoques recientes como el de \citet{eisenstein2023identification} han demostrado que las configuraciones extremas del viento pueden clasificarse en tipos discretos mediante aprendizaje automático, la variabilidad espacial continua de la estructura ciclónica ---particularmente la coherencia entre viento y precipitación--- requiere herramientas de reducción dimensional que preserven la interpretabilidad física. En este sentido, el análisis de Funciones Ortogonales Empíricas Multivariadas (MEOF) se establece como el marco apropiado \citep{hannachi2007empirical}: A diferencia del análisis univariado de EOF aplicado separadamente a cada campo, la técnica MEOF permite extraer modos de variabilidad inherentemente acoplados entre variables físicamente distintas \citep{liang2018multivariate}, revelando la organización espacial coherente donde la circulación del viento y los núcleos de precipitación covarían durante la evolución del sistema.
%
% Liang (2018) Abstract: "The MEOF approach allows for the extraction of coupled patterns of variability shared among the variables... constructed from nitrate, salinity, and potential temperature..." (Valida el uso de MEOF para encontrar patrones compartidos entre variables distintas).
% TEXTO LITERAL LIANG ET AL. (2018) p. 1506: "The multivariate EOF (MEOF) method, a variant of the archetypal EOF method, has been widely used for investigating large-scale atmospheric and oceanic coupled variability structures because of its salient ability to incorporate different variables with their combined variances."
%
La aplicación de este marco estadístico sobre las fases dinámicas previamente clasificadas permitirá construir "prototipos estructurales", sintetizando la estructura típica del ciclón para cada región de génesis y estadio evolutivo, superando así la visión estática de los promedios climatológicos.
% ---% ---% ---% ---% ---% ---% ---% ---% ---% ---% ---% ---% ---% ---% ---% ---% ---% ---% ---% ---
% --- PÁRRAFO 8: DATOS Y HERRAMIENTAS (ERA5 Y TRACKING POR VORTICIDAD) ---

La viabilidad de esta caracterización estructural reside en la reciente disponibilidad del reanálisis ERA5 \citep{hersbach2020era5}, cuya alta resolución espacial ($\sim$31 km) y temporal (horaria) permite resolver los gradientes de presión y los flujos de humedad que generaciones anteriores de datos suavizaban. Validaciones específicas para el Atlántico Sur realizadas por \citet{gramcianinov2020analysis} confirman que ERA5 ofrece una representación superior de la intensidad de los vientos superficiales y la estructura de los ciclones en comparación con productos previos como CFSR.
% Gramcianinov (2020): Validates ERA5 superiority in the South Atlantic.
No obstante, variables de ERA5 en superficie asociadas a ciclones demuestran que, si bien el producto representa bien el comportamiento promedio de los sistemas, tiende a subestimar la magnitud de los extremos más intensos \citep{chen2024evaluation}.

Para capitalizar esta precisión, este estudio utiliza la base de datos de trayectorias de \citet{coutodesouza2024new}, la cual se fundamenta en el criterio de la vorticidad relativa en 850 hPa ($\zeta_{850}$). Esta elección metodológica posee un consenso transversal en la literatura: para resolver el problema de aislar la circulación del sistema, la recomendación de \citet{sinclair1997objective} es utilizar $\zeta_{850}$ como variable de seguimiento, una práctica establecida clásicamente por \citet{sinclair1994objective} y consolidada por \citet{padilhareinke2024objective} como el estándar más robusto para el Hemisferio Sur. La literatura reciente utiliza $\zeta_{850}$ para identificar con precisión las regiones de ciclogénesis \citep{gramcianinov2019properties}, mientras otros estudios de validación demuestran su superioridad sobre campos de presión \citep{padilhareinke2024objective}, siendo incluso utilizado en estudios recientes del Hemisferio Norte \citep{chen2025characteristics}. 
% Couto de Souza (2024): Database provider.
% Sinclair (1994, 1997) & Padilha Reinke (2024): Theoretical justification (Classic + SH context).
% Gramcianinov (2019): Regional application.
% Padilha Reinke (2024) Chen (2025): Modern validation (NH context), showing global methodology consensus.
A diferencia de los algoritmos basados exclusivamente en la presión mínima, el tracking por vorticidad permite el análisis del ciclo de vida de los eventos, capturando incluso aquellos sistemas rápidos que los algoritmos isobáricos tienden a omitir. Esta combinación de datos de alta fidelidad con una detección dinámica sensible constituye la base técnica necesaria para alimentar el análisis multivariado (MEOF) propuesto, permitiendo cuantificar finalmente la evolución de la huella espacial de los eventos compuestos en la región.
% ---% ---% ---% ---% ---% ---% ---% ---% ---% ---% ---% ---% ---% ---% ---% ---% ---% ---% ---% ---
\section{Objetivos del Proyecto de Investigación}
\subsection{Objetivo General}
% Caracterizar la estructura interna y la covariabilidad espacial de los campos de viento y precipitación en los ciclones extratropicales del Atlántico Sur a lo largo de su ciclo de vida, integrando la clasificación por fases evolutivas con análisis multivariado para establecer prototipos estructurales diferenciados por región de génesis.
%
% Caracterizar la evolución de la estructura espacial de los campos de viento y precipitación a través de las fases del ciclo de vida de ciclones extratropicales en el Atlántico Sur, cuantificando sus patrones de variabilidad y asimetría de máximos mediante análisis lagrangiano centrado en el sistema, con estratificación por región de génesis.
%
Caracterizar la evolución de la estructura espacial del viento a 10 metros y precipitación a través de las fases ciclo de vida de ciclones extratropicales en el Atlántico Sur, caracterizando sus patrones de variabilidad y estimación de densidad probabilística de máximos, con estratificación por región de génesis.
%
\subsection{Objetivos Específicos}
\begin{enumerate}
    \item Consolidar una base de datos lagrangiana que asocie los campos horarios de alta resolución de viento y precipitación
    (ERA5) con las trayectorias y fases dinámicas (incipiente, intensificación, madurez y decaimiento) de ciclones
    extratropicales del Atlántico Sur, estratificando la muestra por región de génesis (SBR, LPB y ARG) y por intensidad
    dinámica (Población Global y subconjunto p90).
    \item Cuantificar la distribución estadística y la asimetría espacial de los máximos de viento a 10 m y precipitación
    mediante estimación de densidad por núcleos, caracterizando tanto la distribución de probabilidad de sus magnitudes (KDE
    unidimensional) como la localización preferencial de dichos máximos respecto al centro del ciclón (KDE bidimensional).
    % \item Identificar los modos dominantes de covariabilidad espacial entre los campos de viento y precipitación aplicando
    % Funciones Ortogonales Empíricas Multivariadas (MEOF) en un marco lagrangiano centrado en el sistema, evaluando la
    % coherencia estructural entre ambas variables y su transformación a lo largo del ciclo de vida ciclónico.
    \item Identificar los modos dominantes de variabilidad espacial de los campos de viento y precipitación mediante Funciones
    Ortogonales Empíricas univariadas (EOF) y multivariada (MEOF), evaluando el patron estructural individual de cada variable
    como la covariabilidad acoplada.
    % \item Construir prototipos estructurales para cada combinación de región de génesis y fase evolutiva, integrando
    % los patrones espaciales dominantes del MEOF y las distribuciones probabilísticas de localización de máximos (KDE) en
    % representaciones compuestas que describan la transformación de la organización de viento y precipitación.
    \item Construir prototipos estructurales sintéticos por fase evolutiva y región de génesis, integrando los campos medios
    condicionados al modo dominante de covariabilidad (MEOF) con las densidades probabilísticas de localización de máximos
    (KDE), para cuantificar la magnitud característica de viento y precipitación.
    

\end{enumerate}
% \begin{enumerate}
%     \item Consolidar una base de datos lagrangiana que asocie los campos de superficie de alta resolución (ERA5) con las trayectorias y fases dinámicas (incipiente, intensificación, madurez, decaimiento) identificadas para las tres regiones de ciclogénesis (SBR, LPB, ARG).
%     \item Determinar la asimetría espacial de los máximos locales mediante la estimación de densidad por núcleos (KDE), cuantificando la probabilidad de ocurrencia de máximos de viento a 10 metros y precipitación respecto al centro del ciclón para identificar el desacoplamiento físico entre la presión mínima y la severidad en superficie.
%     \item Identificar los modos de variabilidad acoplada aplicando Funciones Ortogonales Empíricas Multivariadas (MEOF) a los campos conjuntos de viento y precipitación, para revelar las estructuras coherentes de mesoescala que gobiernan los eventos compuestos en cada fase evolutiva.
%     \item Construir modelos conceptuales (prototipos) para cada región de génesis, sintetizando los patrones dominantes obtenidos (MEOF) y las distribuciones de probabilidad de máximos (KDE) en una representación compuesta que describa la evolución típica de la severidad desde la génesis hasta la lisis.
% \end{enumerate}